\section{Open questions and next steps}

The following questions remain open after this replication. Each has a concrete, free-endpoint-only follow-up.

\begin{enumerate}
  \item \textbf{Coherent + readout errors vs.\ pure depolarizing.} Does ZNE / CDR still measurably reduce error when the noise includes coherent (unitary) errors and asymmetric readout, rather than idealized depolarizing? Our replication used only Qiskit-Aer depolarizing channels; the paper's Fig.~3 uses real hardware traces. \emph{Next:} rebuild the ZNE sweep with \texttt{NoiseModel.from\_backend(FakeManilaV2())} plus a deterministic $\sim$2\% $R_z$ over-rotation on every 1q gate; report $\Delta$-error-reduction vs.\ coherent-error amplitude for each inference method.

  \item \textbf{CDR variants (vnCDR, learning-based).} How does base CDR compare to the vnCDR and learning-based CDR variants that Mitiq now exposes but that post-date the 2022 paper? \emph{Next:} add \texttt{code/cdr\_variants.py} comparing base CDR, vnCDR (with 2--3 folded noise scale factors), and a learning-based CDR on the same 7-op Cirq circuit; sweep \texttt{fraction\_non\_clifford} $\in\{0.1, 0.2, 0.4, 0.6\}$.

  \item \textbf{ZNE + PEC composition.} Would combining ZNE with PEC via Mitiq's composable executor pattern beat either alone? PEC was not exercised in this replication. \emph{Next:} write \texttt{code/pec\_replicate.py} using \texttt{represent\_operation\_with\_local\_depolarizing\_noise}, then \texttt{code/zne\_pec\_combined.py} wrapping the PEC-mitigated executor inside \texttt{execute\_with\_zne}. Report mean $|\text{error}|$ for \{raw, ZNE, PEC, ZNE+PEC\} on 10 seeds.

  \item \textbf{Depth scaling of ZNE advantage.} Our replication used a single circuit depth (8). Does the $\ge 40\%$ error reduction hold at depths 16, 32, 64, or does mitigation degrade as the raw signal approaches the $0.25$ mixed-state floor? \emph{Next:} parametrize \texttt{code/zne\_replicate.py} by depth $\in \{4,8,16,32,64\}$; plot mean absolute error vs.\ depth for each inference method; identify the depth at which each method loses its mitigation advantage.

  \item \textbf{Modern NISQ hardware.} Can the mitigation advantage be reproduced on 2026-era hardware (IBM Heron, IonQ Forte, Quantinuum H2) whose native 2q gate errors are $5$--$10\times$ lower than 2021 IBMQ London / Rigetti Aspen-8? Does ZNE still deliver $\sim$40\% error reduction, or does the mitigation ceiling shrink as hardware improves? \emph{Next:} run \texttt{code/zne\_replicate.py} through \texttt{qiskit\_ibm\_runtime.SamplerV2} on an \texttt{ibm\_kingston} or \texttt{ibm\_torino} backend via the IBM Quantum Open Plan (free tier; $\le$30 total circuits). Fallback: \texttt{NoiseModel.from\_backend(FakeKingston())}. Do \emph{not} use paid QPUs.
\end{enumerate}
